\section{Conclusiones}
La implementación de este pipeline estructurado ha logrado sistematizar el preprocesamiento de grandes volúmenes de datos numéricos. Se comprobó que el análisis de correlación de Pearson es altamente efectivo para la selección de características lineales puras, mientras que métodos como PCA permiten una compresión agresiva del espacio dimensional conservando una alta densidad de información (hasta $98.90\%$ en el dataset Vh).

Contar con datasets rigurosamente normalizados, purgados de variables redundantes y estratificados metodológicamente garantiza que los modelos de clasificación posteriores operen con la máxima eficiencia espacial y computacional posible, previniendo sobreajustes y asegurando la validez del modelo final.